%==============================================================================
% RESULTS - PDB Validation Section (NEW)
%==============================================================================
\section{Results}

\subsection{PDB Structure Validation}

To independently validate the CCD mapper accuracy, we compared GlycoSMILES2BAP output against known correct structures from the Protein Data Bank. This validation approach was necessitated by current restrictions on AF3 model parameter access, providing an alternative but equally rigorous validation pathway.

\subsubsection{Validation Dataset}

Four PDB structures were selected to test critical stereochemistry handling capabilities:

\begin{table}[h]
\centering
\caption{PDB Validation Test Cases}
\begin{tabular}{lllp{5cm}}
\toprule
ID & PDB & Structure & Key Validation Point \\
\midrule
V001 & 5NSC & Fucosylated & FUC = alpha-L-fucose (NOT beta) \\
V002 & 1L6R & Lactose & GAL C4 axial hydroxyl (epimer distinction) \\
V003 & 2VXR & Sialyllactose & SIA anomeric = C2 (ketose handling) \\
V004 & 5K65 & M3 N-glycan & Branch topology preservation \\
\bottomrule
\end{tabular}
\label{tab:pdb_validation}
\end{table}

\subsubsection{Validation Results}

GlycoSMILES2BAP achieved 100\% accuracy on all four test cases:

\begin{table}[h]
\centering
\caption{PDB Validation Results Summary}
\begin{tabular}{llcccc}
\toprule
Test Case & PDB ID & Expected CCD & Actual CCD & Anomeric & Status \\
\midrule
Fucosylated & 5NSC & FUC & FUC & C1 & PASS \\
Lactose & 1L6R & GAL & GAL & C1 & PASS \\
Sialyllactose & 2VXR & SIA & SIA & C2 & PASS \\
M3 N-glycan & 5K65 & MAN/BMA/NAG & MAN/BMA/NAG & C1 & PASS \\
\bottomrule
\end{tabular}
\label{tab:pdb_results}
\end{table}

\textbf{Critical Finding - Sialic Acid C2 Anomeric Position:}

The V003 (Sialyllactose, PDB:2VXR) test case represents the most critical validation. Sialic acids (Neu5Ac) are ketoses with the anomeric position at C2, not C1. The CCD mapper correctly:
\begin{enumerate}
\item Maps Neu5Ac to SIA (not alternative codes)
\item Assigns anomeric carbon to C2 (validated against PDB structure)
\item Assigns ring oxygen to O6 (6-membered ring, not O5)
\end{enumerate}

This validation directly addresses the stereochemistry error pattern identified by Huang et al. (2025), where SMILES-based approaches incorrectly place sialic acid anomeric bonds at C1.

\subsubsection{Detailed Validation by Structure}

\textbf{V001: Fucosylated Structure (PDB:5NSC)}

\begin{itemize}
\item Input: Fuc(a1-2)Gal(b1-4)Glc
\item Expected CCD codes: [FUC, GAL, GLC]
\item Actual CCD codes: [FUC, GAL, GLC]
\item Fucose correctly mapped to FUC (alpha-L), not BDF (beta-L)
\item L-configuration preserved through CCD lookup
\end{itemize}

\textbf{V002: Lactose (PDB:1L6R)}

\begin{itemize}
\item Input: Gal(b1-4)Glc
\item Expected CCD codes: [GAL, GLC]
\item Actual CCD codes: [GAL, GLC]
\item Galactose correctly distinguished from glucose (C4 axial vs equatorial)
\item Beta linkage preserved
\end{itemize}

\textbf{V003: Sialyllactose (PDB:2VXR)}

\begin{itemize}
\item Input: Neu5Ac(a2-3)Gal(b1-4)Glc
\item Expected CCD codes: [SIA, GAL, GLC]
\item Actual CCD codes: [SIA, GAL, GLC]
\item SIA anomeric position: C2 (CORRECT for ketose)
\item SIA ring oxygen: O6 (6-membered ring)
\item Alpha configuration preserved
\end